\documentclass[a4paper,12pt]{article}
\usepackage[finnish]{babel}
\usepackage[utf8]{inputenc}
\usepackage[T1]{fontenc}
\usepackage{lmodern}
\usepackage{eurosym}
\usepackage[pdftex,colorlinks]{hyperref}
\title{Etelä-Espoon sosialidemokraattinen työväenyhdistys - toimintakertomus vuodelta 2025}
\author{EED}
\date{22.03.2026}
\begin{document}
\maketitle
\section{Yleistä}
Vuonna 2025 pidettiin kunta- ja aluevaalit, jotka sujuivat SDP:n osalta hyvin Espoossa ja Länsi-Uudellamaalla. Kaksi yhdistyksemme jäsentä, Markku Sistonen ja Helena Tiitinen tulivat valituksi Espoon kaupunginvaltuustoon. Helena Tiitinen valittiin lisäksi Uudenmaan hyvinvointialueen valtuustoon. Yhdistyksen johtokunta kokoontui säännöllisesti.
\section{Vuoden 2025 tavoitteiden toteutuminen}
Yhdistyksen tavoitteet kunta- ja aluevaalien osalta toteutuivat hyvin. Puolueemme presidenttiehdokas. Jäsenmäärässä ei tapahtunut merkittävää kasvua. Verkkoviestintää ei saatu tehostettua tarpeeksi ja sivustomme uudistaminen on kesken. Yhteinen matka Kyprokselle toteutui. Yhdistyksen jäsenet toimivat aktiivisesti Työväen Sivistysliiton Espoon ja Kauniaisten opintojärjestön hallituksessa.
\section{Talous}
Yhdistyksen taloudellinen tilanne pysyi hyvänä koko vuoden ja merkittävä osa omaisuudesta on sijoitusrahastoissa sekä Falklammen kiinteistössä. Suunnitellun mukaan yhdistys luopui tavoitteestaan hankkia uutta sijoitusasuntoa k.o. omaisuuden työlään ylläpidon vuoksi.
\section{Tapahtumien järjestäminen}
Yhdistys osallistui kunta- ja aluevaalien kampanjointiin, järjesti laskiaistapahtuman ja ruusujen jakoa kauppakeskuksissa.
\section{Poliittiset kannanotot}
Yhdistys ei tehnyt omissa nimissään merkittäviä poliittisia kannanottoja hallituksen toimintaan. Puolue on kuitenkin ollut asiassa hyvin aktiivinen.
\section{Luottamustehtävät}
Yhdistyksemme jäsenillä oli seuraavia merkittäviä luottamustehtäviä vuonna 2024 valtiollisissa hallintoelimissä, Espoon kaupungilla, Uudenmaan hyvinvointialueella ja puolueemme järjestöissä. HUOM: Osa tehtävistä on muuttunut syyskokouksen 2024 valintojen vuoksi tälle vuodelle!
\begin{itemize}
\item{Europarlamentti:} Maria Guzenina
\item{Eduskunta:} Miapetra Kumpula-Natri
\item{Aluevaltuusto:} Helena Tiitinen
\item{Kaupunginvaltuusto:} Markku Sistonen (puheenjohtaja), Helena Tiitinen
\item{Lautakunnat:} Markku Sistonen, HUS:in hallituksen varapuheenjohtaja, Leo Heiskanen, liikuntalautakunnan puheenjohtaja, Helena Tiitinen (liikuntalautakunta), Jonne Hänninen (lautamies), Martti Hellstöm (kestävä Espoo, kasvun ja oppimisen lautakunta)
\item{Uudenmaan piirijärjestön edustajisto:} Jonne Hänninen, Mikko Nummelin (varalla), Helena Tiitinen, Sinikka Suontio (varalla)
\item{Kunnallisjärjestön edustajisto:} Raimo Helén, Markku Sistonen, Pauliina Sistonen, Sinikka Suontio. Varalla sisääntulojärjestyksessä: Mikko Nummelin, Liisa Kivimäki, Niina Hänninen, Janne Koski.
\item{Kunnallisjärjestön hallitus:} Mikko Nummelin, Raimo Helén (seuraajajäsen)
\item{Naispiirijärjestön edustajat:} Anja Jaatinen, Anita Paatelma, Helena Tiitinen.
\item{Seurakuntavaltuutetut:} Martti Hellström (tuomiokirkkoseurakunta), Anita Paatelma (Espoonlahden seurakunta), Mikko Nummelin (Olarin seurakunta, varalla)
\end{itemize}
\end{document}
